\documentclass[11pt,a4paper]{article}
\usepackage[margin=24mm,headheight=14pt]{geometry}
\usepackage[T1]{fontenc}
\usepackage{lmodern,microtype,amsmath,amssymb,amsthm,mathtools}
\usepackage{booktabs,tabularx,array,enumitem}
\usepackage[dvipsnames]{xcolor}
\usepackage{xurl}
\usepackage[colorlinks=true,linkcolor=MidnightBlue,urlcolor=MidnightBlue]{hyperref}
\usepackage{fancyhdr}
\pagestyle{fancy}\fancyhf{}
\fancyhead[L]{\small BundleFuse: proposed experiments}
\fancyhead[R]{\small Theory, accuracy and cost}
\fancyfoot[C]{\thepage}
\renewcommand{\headrulewidth}{0.3pt}
\setlength{\parindent}{0pt}
\setlength{\parskip}{6pt plus 1pt}
\setlength{\emergencystretch}{2em}
\clubpenalty=10000\widowpenalty=10000
\setlist[itemize]{leftmargin=1.5em,itemsep=2pt,topsep=3pt}
\setlist[enumerate]{leftmargin=1.7em,itemsep=2pt,topsep=3pt}
\newcolumntype{Y}{>{\raggedright\arraybackslash}X}
\newcommand{\E}{\mathbb E}
\newcommand{\ind}{\mathbf1}
\newcommand{\norm}[1]{\left\lVert#1\right\rVert}
\newcommand{\rownorm}{\operatorname{rownorm}}
\newtheorem{observation}{Observation}
\hypersetup{pdftitle={Experiments to Test and Improve BundleFuse},pdfauthor={BundleFuse research analysis}}

\begin{document}
\thispagestyle{plain}
{\Large\bfseries Experiments to Test and Improve BundleFuse\par}
{\large A prioritized plan following the score-sufficiency analysis\par}
{\small Proposed protocol \hfill 7 November 2026\par}
\vspace{3pt}\hrule\vspace{7pt}

\textbf{Recommendation.} Run four focused studies: remove unnecessary score
columns; separate individual-item support from joint support; test whether
query-specific support improves the frequency fallback; and test an additional
order of historical interaction on Spotify. Reserve bridge-specific and
learned diagnostic studies for questions left unresolved by these results.

\textbf{Status.} This document proposes experiments. No new model fitting or
benchmark evaluation has been launched for it. The query counts and storage
estimates below were calculated from retained artifacts. All existing
checkpoints, score caches and completed results are to be preserved.

\section{What the next experiments need to establish}

The sufficiency analysis separates three questions:
\begin{enumerate}
\item \textbf{Necessity:} can a smaller selected score set retain the ranking
quality of the current estimator?
\item \textbf{Estimation:} do the current weighting and fallback rules use
the available relationships effectively?
\item \textbf{Missing information:} is there a useful relationship that the
current score set cannot represent?
\end{enumerate}
These require different controls. A component-removal experiment addresses
necessity. A new combiner can improve estimation without adding information.
A higher-order count can add a missing relationship, but only on queries
large enough for that relationship to exist.

The earlier studies already compared equal and tuned weighted sums,
rank fusion, two learned fusion procedures, frequency fallback and amplified
historical neighbors. The five-dataset follow-up also included constant-mixture
controls. Repeating the same grids is not the next priority. Their limitations
remain relevant: the learned procedures used limited fitting grids, the source
inventories had been selected earlier, and inspected test sets cannot become
untouched confirmation through another random partition.

\begin{center}
\small
\begin{tabularx}{\linewidth}{@{}p{0.21\linewidth}Yp{0.20\linewidth}@{}}
\toprule
Priority / study & Main question & Initial datasets\\
\midrule
1. E1: compact score set & Which columns can be removed with acceptable loss and a measured cost saving? & POG\_DENSE, POG\\
2. E2: separate history terms & Does coupling individual and joint support hide a better estimator? & POG\_DENSE, POG\\
3. E3: support-aware fallback & Does support for this query improve on total candidate frequency? & Fashion first; amplified Spotify next\\
4. E4: interaction order & Does support from three observed items improve rankings? & Spotify\\
Conditional: E5 & Is the bridge coupling necessary, or is existing information combined poorly? & Where E1--E4 leave a question\\
\bottomrule
\end{tabularx}
\end{center}

\newpage
\section{A feasibility finding that changes the priorities}
\label{sec:feasibility}

The saved full-query records give the following counts. These are query
counts, not independent statistical sample sizes; repeated complete bundles
must still be grouped for analysis.

\begin{center}
\begin{tabular}{@{}lrrr@{}}
\toprule
Dataset & Test queries & Queries with at least 3 seeds & Fraction\\
\midrule
POG\_DENSE & 5,938 & 38 & 0.640\%\\
POG & 4,000 & 0 & 0\%\\
Spotify & 4,000 & 3,975 & 99.375\%\\
\bottomrule
\end{tabular}
\end{center}

Here a seed is an observed item. Dense POG has 11 one-seed queries, 5,889
two-seed queries, 36 three-seed queries and two four-seed queries. POG has
1,990 one-seed and 2,010 two-seed queries.

\begin{observation}[An exact ceiling on query interaction order]
Fix the training data, attributes and candidate $c$. Any real-valued set
score $\psi_c(S)$ restricted to sets with $|S|\leq2$ can be written as
\begin{equation}
 \psi_c(S)=u_c+\sum_{s\in S}v_{cs}
                 +\sum_{\{s,t\}\subseteq S}w_{cst}.
 \label{eq:orderceiling}
\end{equation}
\end{observation}
\begin{proof}
Choose $u_c=\psi_c(\varnothing)$, or choose any empty-set extension if the
score is defined only for nonempty queries. Set
$v_{cs}=\psi_c(\{s\})-u_c$ and
$w_{cst}=\psi_c(\{s,t\})-\psi_c(\{s\})-\psi_c(\{t\})+u_c$.
Substitution verifies the three possible query sizes.
\end{proof}

This is a statement about possible query interactions, not a proof that
BundleFuse estimates all their coefficients correctly. Its coefficients are
strongly constrained by counts, features and a small number of scalar choices.
A new method can improve those coefficients even on one- or two-seed queries.
It could also use richer training structure to estimate them. What cannot act
on those queries is a new term requiring three distinct observed items.

Thus the four-item counterexample in the sufficiency analysis is a genuine
universal limitation, but it is not a compelling explanation of errors on
almost all supplied fashion queries. With all other scores and coefficients
fixed, adding a term that is zero for $|S|<3$ can change no POG rankings and
at most 38 dense POG rankings. Its aggregate improvement in either bounded
metric is at most $38/5938\simeq0.0064$. Retuning other coefficients can
change more queries; such changes must not be attributed to the new term.

This observation motivates E2 for fashion and E4 for Spotify. It supplies
a defensible ceiling on one kind of complexity, while leaving the number
and quality of score representations to be tested.

\section{Common notation and experimental controls}
\label{sec:protocol}

Let $B\in\{0,1\}^{m\times n}$ be the training bundle--item matrix, $S$
the observed set, $G$ the hidden set, $c\notin S$ an eligible candidate,
$f_c=\sum_bB_{bc}$ its training frequency, and $d_b=\sum_iB_{bi}$ the
size of training bundle $b$. Define
\begin{equation}
 a_i=(f_i+5)^{-1/2},\quad v_b=(d_b+1)^{-1/2},\quad
 o_b(S)=\sum_{s\in S}B_{bs}a_s,\quad k_b(S)=\sum_{s\in S}B_{bs}.
\end{equation}
The direct score and incidence correction are
\begin{align}
 p_c(S)&=\sum_bB_{bc}v_b o_b(S),\notag\\
 \delta_c(S)&=1.5\sum_bB_{bc}v_b o_b(S)[k_b(S)-1]_+,\notag\\
 g_c(S)&=p_c(S)+\delta_c(S),\qquad [x]_+=\max(x,0).
 \label{eq:history}
\end{align}
For a raw vector $r$, $z(r)$ means the existing per-query catalogue
standardization. The complete selected score is $Y_c$, the scaled restricted
fallback is $F_c=sP_c$, and the existing interpolation is
\begin{equation}
 \theta_c=F_c+\rho_f(c)(Y_c-F_c),\qquad
 \rho_f(c)=\frac{f_c}{f_c+\kappa_f}.
 \label{eq:gate}
\end{equation}
When a gate strength is zero, define its coefficient as one, including at
zero support, matching the existing no-fallback convention.

All primary studies retain full-catalogue ranking, observed-item exclusion,
the same deterministic tie rule and the same metric definitions. The
fitting target is NDCG@20; Recall@20 is a prespecified guard against losing
retrieval quality. Every candidate within a comparison receives the same
training data, feature inputs and grouped fitting/tuning partition.

Use three distinct roles for data:
\begin{itemize}
\item \textbf{Operator construction:} training bundles and permitted
catalogue attributes. No completion target from the fitting or tuning queries
may remain as a training bundle used to construct their historical score.
Audit exact duplicate complete bundles across those roles.
\item \textbf{Development:} fit scalar choices on fitting groups and select
configurations on tuning groups. Keep every query from one complete bundle
$S\cup G$ in one group. Reuse the existing grouped 70/30 validation split
for the first exploratory stage. The previously selected source inventory
still makes this a conditional study, not independent source discovery.
\item \textbf{Confirmation:} use genuinely unexamined complete bundles
after systems, contrasts and thresholds are frozen. Newly obtained data or
a demonstrably unused corpus portion may provide this. A new mask or a
random split of already studied bundles does not by itself do so.
\end{itemize}

The original all-validation checkpoint remains a descriptive reference.
Matched controls are fitted under the new study's common procedure; do not
present a comparison between differently selected references as an isolated
mechanism effect. Development cross-validation or repeated partitions can
measure selection stability, but folds and seeds are not extra independent
test datasets.

\section{E1: identify a compact score set and measure its actual cost}
\label{sec:e1}

\textbf{Hypothesis.} The full selected bank contains replaceable columns,
and some can be removed without a practically important ranking loss.
This directly addresses why that many scores are retained.

\textbf{Initial search.} The frozen dense bank has five selected columns:
category, bridge, popularity, incidence-enhanced history and PPMI-SVD.
Enumerate its $2^5-1=31$ nonempty subsets. The POG bank has four selected
columns: category, style, bridge and direct history. Enumerate its 15
nonempty subsets. For this experiment the effective historical column is
fixed; E2 changes its construction separately.

Run two different removal controls:
\begin{enumerate}
\item \textbf{Fixed-coefficient removal:} remove a source from both $Y$
and $P$ wherever used, without changing the other coefficients or fallback
scale. This measures reliance on that source in the fitted recipe.
\item \textbf{Refitted subset:} fit each legal subset using the same
deterministic coordinate search and gate/scale grids as its full-bank control.
This measures whether other retained sources can replace it.
\end{enumerate}

Removed columns must not reappear through a default fallback list. Define
the fallback from retained designated sources only. If none remain, disable
the fallback with $\kappa_f=0$. Record that structural change explicitly
in the fixed-coefficient removal analysis. Internal dependencies remain:
removing a standalone PPMI score does not remove PPMI needed for SVD, and
keeping the bridge still requires its historical propagation operator.

Use fixed initializations and the same search budget. Retain legal solutions
from nested subsets as additional initializations when fitting supersets,
so a poor local search does not manufacture a benefit from deletion.
Enumeration of subsets is exhaustive, but the coordinate weight search is
not a proof of global optimality. Any minimality claim is conditional on
this bank and fitting procedure. Candidate configurations and training
budgets must be saved, including failed or discarded alternatives.

\textbf{Selection rule.} Proposed practical margins are a loss of at most
1\% of the matched full model's development NDCG and 1\% of its development
Recall. Freeze the resulting absolute margins before confirmation. These
are chosen utility tolerances, not constants derived from the theory.
For example, reference NDCG $0.28$ gives a margin $0.0028$.
Prefer the least-cost qualifying subset. Require a measured improvement of
at least 20\% in warm inference latency or required persistent storage for
an operational efficiency claim; this is a proposed engineering threshold.
Do not loosen either threshold after seeing confirmation outcomes.

\textbf{Cost measurements.} Time exact scoring, standardization, masking
and top-20 retrieval on the same hardware and thread settings. Report batch
one and one fixed batched setting, median and 95th-percentile query latency,
and separate loading/building times. Use a fixed representative query sample,
warm-up passes and repeated measurements. Record peak resident memory,
required uncompressed arrays and serialized checkpoint size separately.
Export a new dependency-pruned artifact if needed; preserve the original
checkpoint. Verify prediction parity with the selected subset before timing.

\textbf{Decision.} A qualifying smaller model supports a bounded adequacy
claim and can become the proposed compact version. If deletion consistently
hurts, the experiment supplies evidence for retention. If uncertainty crosses
the tolerance, report insufficient precision. If fewer columns do not reduce
actual cost because their dependencies remain, claim simplification only.

After a successful pilot, Spotify's eight-column bank has 255 nonempty
subsets. Measure replay time before committing to enumeration. A smaller
backward-deletion search is an acceptable exploratory alternative, but its
output must not be called the minimum over all 255 subsets.

\section{E2: separate individual support from joint support}
\label{sec:e2}

\textbf{Hypothesis.} For the one- and two-seed fashion regime, the useful
relationships may already be present, while their current coupling is too
restrictive. A better weighting can improve results without adding a new
interaction order.

Hold the selected nonhistorical columns fixed and compare these historical
contributions, where all outer weights are nonnegative:
\begin{align}
 H^{\rm pair}_c &=w_p z(p)_c,\notag\\
 H^{\rm coupled}_c
 &=w_h\{(1-\alpha)z(p)_c+\alpha z(p+\delta)_c\},\notag\\
 H^{\rm separate}_c &=w_p z(p)_c+w_\delta z(\delta)_c,\notag\\
 H^{\rm joint}_c &=w_\delta z(\delta)_c.
 \label{eq:histcontrols}
\end{align}
The coupled control includes the existing $\alpha$ grid
$\{0,0.5,1\}$. The separate version can set either weight to zero. Use
the established scalar weight grid $\{0,0.25,0.5,1,1.5\}$ initially.
First freeze the other source weights, fallback scale and gate, tuning only
the historical contribution. Then run a clearly labeled refitted comparison
with the same permitted outer fitting procedure for all four controls.

For $S=\{s,t\}$, the joint term is exactly
\begin{equation}
 \delta_c(S)=1.5(a_s+a_t)\sum_bv_bB_{bc}B_{bs}B_{bt}.
\end{equation}
The prefactor is constant across candidates for this query. Thus its
standardized ranking is that of a weighted count of bundles containing the
candidate and \emph{both} observed items, apart from numerical stabilizers.
For one-seed queries, $\delta$ is zero. This produces a direct, testable
interpretation for almost the entire fashion evaluation regime.

Do not derive $z(\delta)$ by subtracting $z(p)$ from $z(g)$: each vector
has a different mean and standard deviation. Compute the raw correction
from training incidence or retained raw arrays, then standardize it once.
Keep the existing historical-bundle amplified comparator as a matched
control, reusing valid arrays where possible. This is a test of coefficient
constraints, not a claim of inventing joint-occurrence scoring.

\textbf{Measurements.} Report full-set metrics, one- versus two-seed
results, frequency strata and the fraction of two-seed queries with any
nonzero joint historical support. Strata are defined from observed inputs
before examining outcomes. A zero joint weight or improvement caused only
by retuning other sources cannot support a joint-support contribution.

\textbf{Decision.} If the separate construction improves the matched
coupled version, replace the restrictive coupling with an explicitly
estimated balance, subject to confirmation. If the coupled version is
adequate, its smaller fitting problem gains empirical support. If a joint-only
score suffices for two-seed queries, investigate a simpler conditional rule;
do not silently apply it to one-seed queries where it is identically zero.

\section{E3: does support for this query improve the fallback?}
\label{sec:e3}

\textbf{Hypothesis.} Total frequency can overstate support: an item may
appear often but rarely in historical bundles resembling $S$. A quantity
computed from contributing bundles may improve the gate on the explicit
historical score. It is a proposed support measure, not a proven noise variance.

Let $h_b(S)\geq0$ be bundle $b$'s contribution before multiplying by
candidate membership. For the incidence score,
\begin{equation}
 h_b(S)=v_bo_b(S)\{1+1.5[k_b(S)-1]_+\}.
\end{equation}
For a pair or amplified bank, use that bank's own nonnegative contribution,
including its retained-neighbor mask. Define
\begin{align}
 A_c(S)&=\sum_b B_{bc}h_b(S),\qquad
 Q_c(S)=\sum_b B_{bc}h_b(S)^2,\notag\\
 n_{\rm eff}(c,S)&=
 \begin{cases}A_c(S)^2/Q_c(S),&Q_c(S)>0,\\0,&Q_c(S)=0.\end{cases}
 \label{eq:ess}
\end{align}
This number is high when support is spread across many bundles and low
when one or a few dominate. By Cauchy--Schwarz it is at most the number
of contributing bundles, which is at most $f_c$. It is unchanged by a
common positive rescaling of all $h_b$ for a query.

Let $H_c$ be the weighted, standardized \emph{explicit historical column}
in $Y_c$. Define $A^{\rm rest}_c=Y_c-F_c-H_c$. Compare the original
gate to the term-specific proposal
\begin{align}
 \theta^{\rm freq}_c
   &=F_c+\rho_f(c)\{A^{\rm rest}_c+H_c\},\notag\\
 \theta^{\rm support}_c
   &=F_c+\rho_f(c)A^{\rm rest}_c
          +\rho_q(c,S)H_c,\notag\\
 \rho_q(c,S)&=\frac{n_{\rm eff}(c,S)}{n_{\rm eff}(c,S)+\kappa_q}.
 \label{eq:supportgate}
\end{align}
Only the explicit historical column gets the new gate. A measure derived
from that column need not describe the reliability of SVD, category or
bridge scores, so replacing the gate on every source would have a different,
less specific interpretation.

\textbf{Required controls.} In the primary comparison, fix $Y,F,H$ and
the original $\rho_f$ wherever it multiplies $A^{\rm rest}$. Compare the
support coefficient on $H$ with (i) the original frequency coefficient,
(ii) a separately tuned coefficient $f_c/(f_c+\kappa_H)$, (iii) a tuned
constant coefficient, and (iv) a coefficient based on the unweighted number
of contributing bundles. Each of the last three controls has one tunable
scalar, as does the support proposal. This separates query-specific support
from merely adding a tuning parameter or attenuating the historical score.
The unweighted-count control tests whether contribution concentration matters.
Use $\{0,0.05,0.25,0.5,0.75,1\}$ for the constant historical coefficient
and for the complete constant-mixture fraction below.

Also retain the complete no-fallback model and the tuned constant mixture
of $Y$ and $F$ as reference controls. Use an identical fallback scale in
the primary comparison. A secondary comparison may give every eligible
method the same two scale options: unit scale and the fitting-set matched
scale. Attribute gains there to the selected gate/scale procedure, not the
gate alone. Start with the existing $\kappa$ grid
$\{0,0.05,0.5,1,5,20\}$; any extension is selected on development only.
If the fitted historical coefficient is zero, this term-specific experiment
cannot change that model and should be recorded as inactive.

\textbf{Measurements.} Report rankings for cold candidates, warm candidates
with zero query overlap, and candidates whose support is concentrated or
distributed. Define the bins on fitting data, freeze their edges, and report
denominators. Keep catalogue standardization fixed: at zero raw support,
a standardized historical score can be negative, so changing its gate also
changes a frequency-dependent offset. Describe that effect rather than
interpreting all gains as reduced statistical variance.

\textbf{Optional mechanism check after a promising result.} Retain nested
25\%, 50\% and 100\% subsets of training bundles using three predeclared
hash seeds, keeping evaluation queries fixed. Rebuild \emph{all} historical
operators, frequencies, category associations, SVD and bridge from each
retained training set. Hold supplied descriptors fixed. Compare gates with
matched fitting at each level, and separately report a frozen-coefficient
stress test. This intervention changes available historical evidence; it
does not prove that frequency is the causal source of every ranking gain.
Repeated training seeds are not additional independent query samples.

\textbf{Decision.} A gain over the separately tuned frequency and constant
historical-term controls, as well as the original model,
would support the proposed query-specific correction. A tie favors the
simpler gate. A failure in the zero-overlap or rare-item strata argues against
using this support measure. The existing negative variance-fit slopes alone
justify neither a normal-noise model nor this new gate.

\section{E4: test the missing interaction order where it can act}
\label{sec:e4}

\textbf{Primary dataset: Spotify.} The supplied POG test queries cannot
activate this term. Dense POG is useful only for an explicitly small eligible
subset analysis, not the primary search for a global improvement.

For an integer $j\geq1$, define
\begin{equation}
 r^{[j]}_c(S)=\sum_b B_{bc}v_bo_b(S)
                         \binom{k_b(S)-1}{j-1},
 \label{eq:orders}
\end{equation}
where the combinatorial factor is defined as zero when $k_b(S)<j$.
Then $r^{[1]}=p$, $1.5r^{[2]}=\delta$, and the proposed additional term is
\begin{equation}
 r^{[3]}_c(S)=\sum_bB_{bc}v_bo_b(S)
                 \frac{[k_b(S)-1]_+[k_b(S)-2]_+}{2}.
 \label{eq:third}
\end{equation}
The product is zero for $k_b<3$. Its expansion is
\begin{equation}
 r^{[3]}_c(S)=\sum_{\{s,t,u\}\subseteq S}(a_s+a_t+a_u)
                   \sum_bv_bB_{bc}B_{bs}B_{bt}B_{bu}.
\end{equation}
Each seed weight occurs once for each choice of two other overlapping
seeds. Thus this term measures four-item occurrences: a candidate together
with three observed items. Compute it through overlaps and incidence
multiplication; never materialize an item-by-item-by-item-by-item tensor.

Add a single nonnegative coefficient $\lambda z(r^{[3]})$ to the
historical side of the selected score, with $\lambda$ initially in
$\{0,0.25,0.5,1,1.5\}$. Its fallback membership is fixed to absent.
First hold every existing coefficient and the frequency gate fixed, then
compare equally refitted versions. Record whether the added term is active,
its independent fitting cost, and its incremental scoring cost.

\textbf{Strong controls are essential.} Compare against the retained
amplified Spotify bank with its matched fallback. Include full-neighbor
power-two and power-three historical overlap, and the existing truncated
neighbor alternative. For equal seed weights and fixed bundle size, cubing
the overlap contains a third-order query term:
\begin{equation}
 k^3=k+3k(k-1)+k(k-1)(k-2).
\end{equation}
Top-neighbor truncation can also introduce dependence on the whole query.
Consequently, four-item sensitivity is not a capability unique to the new
count. The proposed advantage, if any, is estimating its contribution
separately instead of coupling lower and higher orders through a fixed power
or neighbor-selection rule.

\textbf{Synthetic check before real ranking.} Reuse the parity-style
four-group construction from the sufficiency analysis. Verify that pair and
triple-based controls fail to distinguish its completions, while the added
term and an appropriate power-three control distinguish them. Verify that
the new term is exactly zero when $|S|<3$ and that setting $\lambda=0$
reproduces the baseline. These are representation checks, not benchmark gains.

\textbf{Decision.} Retain the term only if it improves the matched strong
control beyond uncertainty with acceptable cost. If it helps only a narrow
query-size stratum, report that scope. If amplified neighbors do equally
well, favor the simpler implementation and describe the shared capability.
Do not continue automatically to orders four, five and beyond: an additional
order needs an unresolved error pattern, enough eligible queries, and a
separate development justification.

\section{Conditional studies: coupling and unresolved prediction errors}
\label{sec:e5}

\subsection{E5a: test the bridge only if E1 retains it}

Let $E$ be the matrix of normalized category-residual description vectors.
Let $L$ be the fixed historical propagation operator and
$N=\rownorm(LE)$. For bridge strength $\eta$, let $D$ be the diagonal
row-normalization matrix for $E+\eta N$. The bridge kernel is
\begin{equation}
 D(EE^\top+\eta EN^\top+\eta NE^\top+\eta^2NN^\top)D.
\end{equation}
The two cross terms compare an item's own description with another item's
historical-neighbor description. Compare the bridge with own-description
similarity, neighbor-description similarity, and a tuned weighted sum of the
two separate scores.

For an algebraically controlled ablation, remove only the two cross terms
while holding $D$ fixed. This isolates their contribution to the numerator;
$D$ still contains the original normalization information, so this is not
complete removal of all cross-term effects. Also evaluate a practical
concatenated representation $\rownorm([E,\eta N])$, which has its own
normalization. Treat these as distinct comparisons. Compute scores through
feature products and query means, not dense catalogue-squared matrices.

A gain over a properly tuned separate-score combination would support the
specific coupling mechanism. If a simpler description score suffices, remove
the extra computation. This study was not established by the earlier
historical-bank comparisons, which kept the bridge fixed.

\subsection{E5b: a learned diagnostic, not a new production network}

The theory defines additional information through the change in the best
prediction of normalized relevance after revealing another score. That
population quantity is not directly observed. A small prediction probe can
find remaining errors, but failure of the probe cannot prove sufficiency.

For NDCG@20, define
\begin{equation}
 V_c=\frac{\ind\{c\in G\}}{D_{20}(G)},\qquad
 D_{20}(G)=\sum_{j=1}^{\min(20,|G|)}\frac1{\log_2(j+1)}.
\end{equation}
Fit regularized linear predictors of $V_c$ from the retained standardized
scores, their frequency-gated versions, log frequency and seed count.
Compare a predictor using these features with an equally tuned predictor
also receiving one proposed score. Fit and select within grouped development
folds, then evaluate on held-out groups. These probes estimate coefficients;
they are deliberately learned diagnostic models, not parameter-free controls.

If there are $d$ retained columns, the stated baseline has $2d+3$ fitted
coefficients in total: $2d$ for scores and gated scores, one for log frequency,
one for seed count, and an intercept. Thus a five-column bank uses 13
coefficients; a single additional scalar score increases this to 14.
The feature definition is
fixed before fitting; the base probe must not silently omit a gate used by
the production model. If a later tested model uses the query-support gate,
add its corresponding transformed column to both probes and report the
updated coefficient count.

Use all candidates in streamed moment calculations where feasible. If
negative candidates are sampled, use a predefined uniform sample with known
inclusion probabilities and appropriate importance weights. Include all
positives. A balanced positive/negative loss without correction estimates
a different target from full-catalogue squared prediction error. Report
ranking metrics as well as the properly weighted error.

An improvement by the base probe suggests that current fusion leaves
information unused. An improvement only after adding a source suggests
unresolved predictive value \emph{within this probe class}. Neither gives
an upper bound on all possible additions. In particular, a learned estimate
must not be substituted into the theoretical information bound as a certified
upper bound. The earlier weak learned-ranker results do not eliminate this
diagnostic question because their objectives and fitting procedures differed.

\section{Confirmation, decisions and permitted claims}
\label{sec:confirmation}

Primary development comparisons should be frozen before each study is
run. Select a compact candidate from E1, one history construction from E2,
one gate from E3 and, if successful, one Spotify extension from E4. Combine
successful changes only in a separate stage with ablations of each change;
separate gains need not add together.

For confirmation, freeze the exact systems, fitting provenance, target
population, resource limits, contrast list, tolerance values and statistical
family before accessing new outcomes. Use full-catalogue query-macro
NDCG@20 and Recall@20 for comparison with existing reports. Resample complete
bundle groups together for paired uncertainty, preserving query multiplicity.
Account for user/time dependence if the data require a larger grouping unit.
Report group counts and query counts separately.

For a compact model, the lower confidence bound on its utility difference
from the reference must exceed the negative prespecified tolerance. For a
superiority claim, the corresponding lower bound must exceed zero after
the declared multiplicity correction. The exact number of contrasts is
generated from the frozen system manifest; reserve that entire family,
including conditional studies that are to be claimed. A nonsignificant
difference is not equivalence. Bootstrap intervals are approximate; disclose
that qualification rather than presenting them as distribution-free proofs.

To make the finite-family sufficiency claim from the theory, evaluate all
frozen members of the declared family, not just the chosen compact model
and the original full one. The original full model need not be the best
member of the family. If $M$ fixed systems are audited on $N$ independent
bundle units using bounded bundle-average utility, a simultaneous
Hoeffding bound has radius
\begin{equation}
 b=\sqrt{\frac{\log(2M/\delta)}{2N}}
\end{equation}
for one metric and failure probability $\delta$. Allocate confidence across
both metrics as well. Choose the least-cost frozen system whose audit utility
is within a prespecified tolerance $\varepsilon$ of the best observed family
member. It is within
$\varepsilon+2b$ in population utility on the simultaneous event. This
certificate concerns the bundle-average target; do not apply the independent-
unit formula to dependent query rows or silently equate bundle-average and
query-macro metrics.

The elementary bound may be too wide to certify the chosen 1\% relative
margin. Estimate paired variability on development data to plan precision,
and use a justified tighter bound if available. If the audit is too small,
report an inconclusive certificate rather than changing the tolerance.
If no genuinely new audit data exist, complete the exploratory studies and
state that confirmation remains unavailable.

\begin{center}
\small
\begin{tabularx}{\linewidth}{@{}YY@{}}
\toprule
Result & Claim it can support\\
\midrule
Smaller bank retains quality with lower measured cost & Some score columns are unnecessary within the declared bank, fitting procedure and tolerance.\\
Separate history terms improve & Their present coupling is restrictive; better results need not require more interaction orders.\\
Query-support gate improves & Candidate history relevant to this query is useful beyond total frequency under the tested controls.\\
Third-order term improves Spotify & Candidate--three-seed occurrences add useful information in the tested setting, subject to amplified-neighbor controls.\\
No detectable gain & No demonstrated benefit at the experiment's resolution; not universal sufficiency.\\
Fresh finite-family certificate succeeds & The selected system is adequate within its declared family, population, metric, tolerance and resource budget.\\
\bottomrule
\end{tabularx}
\end{center}

\section{Execution budget and retained outputs}
\label{sec:execution}

The first studies reuse training operators and verified score caches where
their definitions match. They fit scalar choices rather than neural scoring
networks. E3 introduces at most one new support-gate strength in its fixed
comparison, and E4 one added source coefficient. E5b fits the explicitly
reported linear coefficients. All tuning effort must be counted; zero
gradient-trained scoring parameters does not mean zero fitted choices.

Storage is material. From the retained catalogue and query counts, one
additional float32 score column over validation and test requires approximately
the following, excluding metadata and temporary blocks:
\begin{center}
\begin{tabular}{@{}lrrr@{}}
\toprule
Dataset & Catalogue items & Validation + test queries & GiB per column\\
\midrule
POG\_DENSE & 31,217 & 8,906 & 1.036\\
POG & 48,676 & 6,000 & 1.088\\
Spotify & 254,155 & 6,000 & 5.681\\
\bottomrule
\end{tabular}
\end{center}
The estimate is $4n(N_{\rm val}+N_{\rm test})/2^{30}$ GiB, where
$N_{\rm val}$ and $N_{\rm test}$ are the validation and test query counts.
Generate development columns first, then materialize test/audit columns only
for frozen candidates. Share arrays across fusions. Check available disk and
record a pilot's time per scoring block and per configuration before projecting
wall time. Existing results do not justify a new fixed runtime estimate.

The proposed implementation location is \path{experiments/sufficiency/};
it does not yet exist as a runnable study. Use separate future result and
checkpoint roots, \path{results/sufficiency/} and
\path{checkpoints/sufficiency/}. Preserve existing files and interrupted
partials. Each run should retain:
\begin{itemize}
\item input, source and protocol hashes; grouping manifests and split indices;
\item legal source subsets, fitting grids, all selected/discarded recipes and
explicit fallback membership;
\item completed cache blocks, completion markers and resumable logs;
\item per-query metrics, top-20 predictions, subgroup counts and paired results;
\item dependency-aware latency/memory reports and parity checks for compact exports.
\end{itemize}
Required correctness checks include baseline equality at zero added weight,
zero joint terms on ineligible query sizes, support bounds in
equation~\eqref{eq:ess}, exclusion of removed fallback sources, and identical
ranking masks/tie rules. Do not change sources during a resumable run.

\textbf{First execution package.} E1 and E2 on POG\_DENSE and POG, followed
by E3 on the surviving banks. Promote E4 to a full Spotify run only after
its synthetic checks and a development pilot. E5 and training-thinning
studies require a specific unresolved finding. This order prioritizes the
question of how many scores are needed while preserving a concrete route
to improved accuracy.

\section*{Evidence behind this plan}

The query-size and storage audit is saved in
\path{results/algebraic_diagnostics/sufficiency_experiment_feasibility.json}.
It records the source hashes of each retained query record, identity manifest
and frozen recipe. It is a read-only feasibility audit, not a new evaluation.

Completed comparison evidence is under
\path{results/novelty_attribution/} and
\path{results/amplified_fallback/}. Their protocols are in
\path{experiments/novelty/README.md} and
\path{experiments/amplified_fallback/README.md}. The mathematical analysis
is available as \path{docs/score_sufficiency.tex} and
\path{output/pdf/score_sufficiency.pdf}; its numerical checks are in
\path{results/algebraic_diagnostics/score_sufficiency.json}.
This proposal builds on those records and does not replace their scope
qualifications or relabel exploratory results as confirmation.

\end{document}
